% =======================================================================================
\section{Introduction}
\label{sec:introduction}
% =======================================================================================

\Ac{lisa} will observe millions of \acp{gb} in the Milky Way, quasi-monochromatic sources
between \qtyrange{\RelatedWorkFzeroLow}{\RelatedWorkFzeroHigh}{\milli\hertz}, and resolve thousands of
them \citep{nelemans2001gravitational,korol2017prospects,amaro2017lisa,colpi2024lisa,karnesis2021characterization}. Below a few millihertz
their signals overlap: a source spreads over several frequency bins through the instrument's
orbit and its own chirp, and neighbours share bins, so their parameters must be inferred jointly
(\cref{fig:related-crowding}). The reference method is a global fit, a trans-dimensional \ac{mcmc}
over every resolvable source \citep{littenberg2020global,littenberg2023prototype,katz2024efficient,strub2024global,deng2025modular}.
It is the standard for what the data support, but even a joint chain over one of our windows takes
\qtyrange{\McmcCostJointLo}{\McmcCostJointHi}{\hour}, and it has to run again on every new stretch
of data. Amortised estimators learn the posterior once and
then return it in a fraction of a second, but for \acp{gb} they have so far treated one source,
two ordered sources in a narrow band, or a factorised set of toy signals
\citep{korsakova2024neural,houba2025slotflow,delmond2026neural} (\cref{tab:related-closest}).

We bring the amortised route to the global fit's unit of work: the joint posterior of the
\RelatedWorkSourcesWord{} sources in a frequency window, \RelatedWorkParameters{} parameters under
the full instrument response, for windows anywhere in the band, drawn by conditional flow
matching \citep{lipman2023flow,liu2023flow,albergo2023building,dax2023flow} with no likelihood
evaluation. Three difficulties stand in the way. The posterior is multimodal by symmetry: each
source has four copies under its polarisation and phase, and the sources of a window can be
exchanged. Its widths span orders of magnitude, from the prior for a quiet source to a small
fraction of a bin in frequency for a loud one, and a flow run on the prior's scale and sampled on
a uniform grid cannot resolve the narrow end. And with no reference posterior for most windows, the usual calibration
checks pass flows that are plainly wrong.

\paragraph{Contributions.}
\begin{itemize}
  \item A flat chart of eight coordinates per source in which the signal is linear in one complex
    amplitude per helicity, so a source's four equivalent copies become one mode and the flow
    runs on straight paths (\cref{sec:method-chart}).
  \item Sources as slots in order of frequency, each with a Gaussian frame the network reads off
    the data alone, around the prior's moments, so the flow transports order-one widths for loud
    and quiet sources alike (\cref{sec:method-chart,sec:method-fm}).
  \item A flow-matching recipe for posteriors that are narrow at the end of the path: Gaussian
    likelihoods of the frame and of the endpoint in it, whose reward for narrowing is relative to
    each source's own width, and an analysis of what the sampler resolves with and without the
    frame (\cref{sec:method-fm}).
  \item A one-way transformer whose data stream depends on the observation alone, so it is encoded
    once per window and reused by every draw and every sampler step, trained on batches simulated
    on the device (\cref{sec:method-network}).
  \item A validation protocol checked on known answers: a toy with an exact posterior, where ranks
    of the parameters miss phase, sky and dependence errors that ranks of data-space quantities
    catch, and a ladder of mock posteriors run through the whole evaluation
    (\cref{sec:validation}).
\end{itemize}
We train on two bands, the narrow band below \qty{\PriorsNarrowHighMhz}{\milli\hertz} with windows
of \PriorsNarrowW{} bins and the full band up to \qty{\PriorsFullHighMhz}{\milli\hertz} with windows
of \PriorsFullW{} bins, and compare against joint and per-source \ac{mcmc} chains
(\cref{sec:results}). \Cref{app:related} covers the related work.
